% H9-T13: generic odd complete intersections in abelian varieties.
% Run JAC1964-20261121; /root/h9_height; GPT-6 Astra Pro.

\subsection*{H9-T13. Generic odd complete intersections in an abelian variety}

\begin{theorem}[H9-T13]
Let $A/k$ be an abelian variety of dimension
$N=2m+c-1$, where $m\ge1$ and $c\ge1$, over an algebraically closed
field of arbitrary characteristic. Let $h$ be an ample divisor class.
Choose a smooth successive complete intersection
\[
 X=A\cap H_1\cap\cdots\cap H_{c-1},\qquad [H_i]=h,
 \qquad \dim X=2m,
\]
with the convention $X=A$ when $c=1$. Let $Y/K$ be the generic member
of a genuine Lefschetz pencil of sections of $X$ in the restricted
class $h$, where $K=k(\mathbf P^1)$. Thus
$\dim Y=2m-1$ and $Y$ has total codimension $c$ in $A$.

Accept 032, Corollary 8.4. The normalized signed incidence class
\[
 \frac{(-1)^m}{2}\Delta^*d
       \quad\text{on }P^m(Y)_{\rm prim}
\]
is a rational polarization. The first $c-1$ smooth sections need not
be generic; genericity is required only of the last member through
the stated genuine Lefschetz-pencil hypothesis.

For a starting projective embedding with class $h_0$, taking
$h=d h_0$ with $d\ge2$ supplies such terminal Lefschetz pencils in
every characteristic. In particular the conclusion holds for a
geometric generic complete intersection of $c$ degree-$d$ sections
of an embedded abelian variety, for every $d\ge2$.

If $m=2$, the actual primitive regular algebraic representative
$\operatorname{Ab}^2(Y)_{\rm prim}$ has the same polarization, and
its canonical map to $P^2(Y)_{\rm prim}$ is an isogeny.
\end{theorem}

\paragraph{The fixed and variable projectors.}
Put $n=2m-1$, let $i:X\hookrightarrow A$, $j:Y\hookrightarrow X$,
and $k=i\circ j:Y\hookrightarrow A$. Weak Lefschetz identifies
$H^n(X)$ with $H^n(A)$ and embeds this space in $H^n(Y)$.
On the ambient subspace the middle pairing is
\[
 (u,v)\longmapsto\int_Auvh^c.
\]
It is nondegenerate, because hard Lefschetz on the $N$-dimensional
abelian variety gives the isomorphism
\[
 L_h^c:H^n(A)\xrightarrow{\ \sim\ }H^{n+2c}(A).
\]
Its inverse is algebraic. Let
\[
 r=(L_h^c)^{-1}k_*:H^n(Y)\longrightarrow H^n(A),
 \qquad r k^*=1.
\]
Choose this inverse supported in the indicated degrees, using the
algebraic K\"unneth projectors of $A$. It gives an algebraic
correspondence $r$. Hence
\[
 \Pi_V=1-k^*r
\]
is the algebraic self-adjoint variable projector on $H^n(Y)$.
This is exactly the variable space for the terminal pencil: the
ambient space of that pencil is $j^*H^n(X)=k^*H^n(A)$.
No separate standard-conjecture hypothesis for $X$ is needed for this
projector.

Let $e_A$ be the algebraic primitive projector on $H^n(A)$, whose
image is $\ker L_h^{c+1}$. Put
\[
 F=k^*\ker L_h^{c+1},\qquad \Pi_F=k^*e_A r.
\]
The ambient primitive decomposition is orthogonal for
$\int_Auvh^c$: if $L_h^{c+1}u=0$ and $v=L_hw$, then
$\int_Auvh^c=\int_A(L_h^{c+1}u)w=0$.
Thus $\Pi_F$ is self-adjoint and orthogonal to $\Pi_V$.

We check precisely that these give the full primitive part on $Y$.
The Gysin map
\[
 k_*:H^{n+2}(Y)\longrightarrow H^{n+2+2c}(A)
\]
is injective by dual weak Lefschetz in degree $n-2$ (and the assertion
is immediate when the source is zero, as for $n=1$).
For an ambient class $k^*u$, its product with $h_Y$ pushes to
$h^{c+1}u$. Therefore it is primitive on $Y$ exactly when
$u\in\ker L_h^{c+1}$. A variable class $v\in\ker j_*$ is primitive
as well: $j_*(h_Yv)=h_Xj_*v=0$, and
$j_*:H^{n+2}(Y)\to H^{n+4}(X)$ is injective by the same dual weak
Lefschetz argument. Consequently
\[
 H^n(Y)_{\rm prim}=F\ \mathbin{\perp}\ E,
 \qquad E=\ker j_*.
                                                               \tag{H9.46}
\]
Incidence detection and Saito's transpose compatibility turn these
algebraic projectors into orthogonal nondegenerate factors of
$P^m(Y)$ in the rational category, as in JAC-T4.1 and H9-T4.

\paragraph{The fixed-part curve test with all exponents.}
Let $C_0$ be a smooth pointed parameter curve of genus $b\ge1$, and
pull a normalized family on the fixed factor back to it. After the
parameter degree-one projection its class is
\[
 \xi\in H^1(C_0)\otimes F.
\]
Set $\eta=(1\otimes r)\xi$. This is algebraic, satisfies
$k^*\eta=\xi$, and obeys $h^{c+1}\eta=0$.
The algebraic inverse to Abel pullback lifts it to an algebraic class
\[
 w\in H^1(J(C_0))\otimes H^n(A),
       \qquad h^{c+1}w=0,
\]
on the abelian variety $J(C_0)\times A$ of dimension
$b+2m+c-1$. The class $w$ has codimension $m$.
Let $\theta$ be the canonical principal polarization of $J(C_0)$
and put $H=\theta+h$. Then
\[
 H^{b+c}w=0.
\]
Indeed every term with $h$-exponent at least $c+1$ vanishes, while
every remaining term has $\theta$-exponent at least $b$ and hence
exceeds the top degree of $J(C_0)$ after multiplication by its
degree-one component. This is exactly the primitive exponent
$(b+2m+c-1)-2m+1=b+c$.

In the quadratic form, degree on the two factors permits only the term
with $h$-exponent $c$ and $\theta$-exponent $b-1$. The Abel curve-class
formula therefore gives
\[
 \begin{aligned}
 \int_{J(C_0)\times A}w^2H^{b+c-1}
  &=\binom{b+c-1}{c}
       \int_{J(C_0)\times A}w^2\theta^{b-1}h^c\\
  &=\frac{(b+c-1)!}{c!}
       \int_{C_0\times A}\eta^2h^c\\
  &=\frac{(b+c-1)!}{c!}
       \int_{C_0\times Y}\xi^2.
 \end{aligned}
 \tag{H9.47}
\]
The coefficient is positive. By 032, the right-hand intersection has
sign $(-1)^m$ and is zero only if $\xi=0$. This proves nefness of the
signed incidence divisor on the fixed factor. Its nondegeneracy was
already supplied by the self-adjoint-factor argument, so the
nef-plus-nondegenerate criterion makes it a rational polarization.
Curves of genus zero have constant maps to an abelian variety and
contribute degree zero; they cause no exception to the curve test.

\paragraph{The variable factor and the terminal pencil hypothesis.}
The terminal ambient variety $X$ is smooth projective of even
dimension $2m$, and its pencil is genuinely Lefschetz. Hard Lefschetz
holds for $X$. SGA 7 II, XVIII, 6.6 and 6.7.2, printed pp.~323--327,
therefore give spanning and conjugacy of the vanishing cycles and
absolute irreducibility after every finite-index monodromy restriction.
Their characteristic-two exception is restricted to odd ambient
dimension, so it is absent here.

If the variable incidence factor is nonzero, its normalized family
gives an injective realization into $E_\ell$ for every invertible
prime. Thus $E_\ell\ne0$, every local vanishing cycle is nonzero,
and the hypotheses of H9-T6 are satisfied, exactly as in H9-T6a.
That theorem proves positivity of the variable factor. If it is zero,
there is nothing to prove on that factor. Combining its polarization
with the fixed polarization through (H9.46) proves the main assertion.

SGA 7 II, XVII, Theorem 2.5.1, printed p.~217, says that a degree-$d$
multiple of a projective embedding is Lefschetz for $d\ge2$. Apply
this to the embedding of the already chosen smooth $X$ induced from
the starting embedding of $A$. It supplies the last-pencil hypothesis
in the stated degree-$d$ construction. Successive Bertini sections
supply the earlier smooth chain. No assertion that an arbitrary
nodal morphism has the global monodromy of such a pencil is used.

\paragraph{The actual representative when $m=2$.}
Existence over the possibly imperfect $K$ and injective realizations
of normalized miniversal families are exactly the source-checked
ACMV results used in H9-T10. They make the same projectors act on
$\operatorname{Ab}^2(Y)$ and give the variable injection required
for H9-T6's sign argument. On the fixed factor, a nonconstant
geometric curve has a nonzero pulled-back family class because its
Tate-module map remains nonzero after that injective realization.
Equation (H9.47), now with $m=2$, makes the incidence degree strictly
positive. The argument excluding a positive-dimensional kernel and
the canonical comparison with the incidence quotient are verbatim
the explicit arguments of H9-T10. They prove the stated rational
polarization and isogeny, without any integral purely inseparable
base-change assertion.\hfill$\square$

\paragraph{Why this does not close the arbitrary fixed-axis gate.}
Let $\kappa:W\hookrightarrow X_0$ be the codimension-two axis in
an abelian variety $X_0$ from H9-T7. The odd variety $W\times E$ is
cut out in $X_0\times E$ by two pullbacks of divisors from $X_0$.
These divisor classes have degree zero along every elliptic fibre
$\{x\}\times E$ and are not ample. Degree on that fibre is constant
in a connected algebraic deformation of the line bundle, so one
cannot make them ample while retaining those divisor classes.

There is also an exact obstruction to using the ambient transfer
above on the remaining axis classes. For
$w\in A_W=\ker(\kappa_*:\operatorname{Alg}^{m-1}(W)
\to H^{2m+2}(X_0))$, the elliptic family $w\boxtimes P_E$
on $E\times(W\times E)$ pushes to
\[
 (\kappa_*w)\boxtimes P_E=0
       \quad\text{in the cohomology of }E\times X_0\times E.
\]
Applying ambient algebraic Lefschetz inverses to this zero class
still gives zero. The transfer in (H9.47) was injective on fixed
cohomology; this attempted transfer kills precisely the unresolved
axis summand. A new injective correspondence or a theorem that lifts
these special classes through a degeneration would be required.
Generic positivity does not supply either assertion. The separate
classical characteristic-zero Hodge--Riemann argument for special
fibres is unaffected by this positive-characteristic limitation.

\paragraph{Status and archive scope.}
This is a deduction written in this run from 032, algebraic Lefschetz
operators on abelian varieties, and the explicitly cited terminal-pencil
monodromy theorems. It extends the primitive incidence construction of
cote 44, scans 8--9, to the indicated generic odd complete intersections.
It neither asserts polarization on every special member in positive
characteristic nor upgrades the original H9 homological domain to an
incidence Jacobian. No historical originality or formal verification
is claimed.
